Search and rescue (SAR) operations demand rapid coverage of large areas, yet ground-based teams are constrained by terrain, fatigue, and limited field of view, resulting in prolonged response times during the critical early hours of an incident \cite{goodrich2008}. This paper presents the design, implementation, and validation of a low-cost autonomous SAR drone that detects casualties using onboard computer vision without reliance on cloud connectivity or ground-station processing.

The system employs a Raspberry Pi~5 running a YOLOv8n object detection model \cite{yolov8} exported to TensorFlow Lite \cite{tflite}, achieving real-time inference at 206\,ms per frame (4.8\,FPS). A finite state machine orchestrates the complete mission: autonomous takeoff, lawnmower search over a user-defined polygon, AI-based casualty detection, GPS-referenced target localisation, controlled descent for operator verification, and precision landing. An operator-in-the-loop confirmation stage ensures that no autonomous landing is initiated on a false positive.

The detection model, trained on 366 mixed synthetic and real-world images at native camera resolution ($1456 \times 1088$), achieves a mean average precision (mAP$_{50}$) of 0.995. A progressive five-stage testing methodology---from software-in-the-loop simulation through bench testing to outdoor field trials---validates each subsystem independently before integration. The complete onboard compute platform costs under \pounds{}100, demonstrating that effective autonomous SAR capability is achievable with consumer-grade hardware and open-source software.
